\documentclass{article}
\usepackage{amsfonts,amsthm,amsmath}
\usepackage[mathscr]{euscript}
\usepackage{enumitem}

\setcounter{section}{0}
\renewcommand{\thesection}{\S\arabic{section}}
\renewcommand{\thesubsection}{\arabic{section}}

\newtheorem{thm}{Theorem}
\newenvironment{theorem}[1]{
	\renewcommand{\thethm}{#1}
	\begin{thm}
		}{\end{thm}}

\newtheorem{lma}{Lemma}
\newenvironment{lemma}[1]{
	\renewcommand{\thelma}{#1}
	\begin{lma}
		}{\end{lma}}

\theoremstyle{definition}
\newtheorem{ex}{Exercise}
\newenvironment{exercise}[1]{
	\renewcommand{\theex}{#1}
	\begin{ex}
		}{\end{ex}}

\title{Exercises 8.D}
\author{Connor Mooneyhan}
\date{March 3, 2024}

\begin{document}

\maketitle

\begin{ex}
	Find the characteristic polynomial and the minimal polynomial of the operator $N$ in Example 8.53.
\end{ex}
\begin{proof}
	From Exercise 8.53, we have that $N \in \mathcal{L}(\mathbf{F}^4)$ is defined by \begin{equation*}
		N(z_1, z_2, z_3, z_4) = (0, z_1, z_2, z_3).
	\end{equation*}
	Since the only eigenvalue of $N$ is 0, it follows immediately that the characteristic polynomial of $N$ is \begin{equation*}
		p(z) = z^4.
	\end{equation*}
	Thus the only possibilities for the minimal polynomial are $z, z^2, z^3$, and $z^4$.
	Clearly $T$ itself is not the zero map, so we must start our search with $z^2$.
	The following calculations are simple: \begin{align*}
		N^2(z_1, z_2, z_3, z_4) & = (0, 0, z_1, z_2), \\
		N^3(z_1, z_2, z_3, z_4) & = (0, 0, 0, z_1),   \\
		N^4(z_1, z_2, z_3, z_4) & = (0, 0, 0, 0),
	\end{align*}
	so in this case $N^4$ is in fact the lowest power of $N$ to be zero, and thus $z^4$ is the minimal polynomial of $T$.
\end{proof}

\begin{ex}
	Find the characteristic polynomial and the minimal polynomial of the operator $N$ in Example 8.54.
\end{ex}
\begin{proof}
	From Exercise 8.53, we have that $N \in \mathcal{L}(\mathbf{F}^4)$ is defined by \begin{equation*}
		N(z_1, z_2, z_3, z_4, z_5, z_6) = (0, z_1, z_2, z_3).
	\end{equation*}
\end{proof}

\end{document}
